\chapter{Wstęp i cel pracy}

Kompilator jest narzędziem, z~którego inżynier oprogramowania korzysta przy
każdym uruchomieniu procesu budowania, a~mimo to rzadko bywa przedmiotem jego
uwagi. W~praktyce traktuje się go jak czarną skrzynkę: na wejściu znajduje się
kod źródłowy, na wyjściu plik wykonywalny, a~to, co dzieje się pomiędzy,
pozostaje poza zainteresowaniem programisty. Poziom optymalizacji ustawia się
zwykle raz, przy konfigurowaniu projektu, i~nie wraca się do niego przez cały
cykl życia oprogramowania. Gdy pojawia się potrzeba przyspieszenia programu,
modyfikowana jest przede wszystkim logika algorytmu, znacznie rzadziej rozważa się
wpływ, jaki na wynik końcowy ma sam proces translacji kodu na instrukcje maszynowe.

Pominięcie to jest o~tyle zaskakujące, że znaczna część oprogramowania
działającego bezpośrednio na sprzęcie powstaje w~językach kompilowanych, takich jak język C
\cite{Tanenbaum}.
Sposób, w~jaki kompilator tłumaczy słowa kluczowe tego języka na instrukcje
konkretnej architektury, decyduje o~liczbie wykonanych operacji, o~wykorzystaniu
rejestrów oraz o~tym, czy program skorzysta z~jednostek wektorowych procesora.
Każdy z~tych czynników wpływa bezpośrednio na czas wykonania, a~jednak żaden
z~nich nie wymaga zmiany choćby jednej linii kodu źródłowego.

Znaczenie tego zagadnienia rośnie w~systemach wbudowanych, ponieważ dysponują one 
 ograniczoną mocą obliczeniową, niewielką pamięcią operacyjną
i~pracują często w rozwiązaniach wymagających oszczędzania energii. W~ich przypadku
poprawa wydajności uzyskana bez wymiany sprzętu przekłada się wprost na niższy
koszt systemu, mniejszy pobór energii oraz większy zapas mocy obliczeniowej
na pozostałe zadania. Jeżeli zmiana ustawień kompilacji pozwala skrócić
czas wykonania krytycznej procedury, jest to korzyść osiągnięta bez żadnych
nakładów.

Wpływu ustawień kompilacji na czas wykonania nie da się jednak oszacować bez
pomiaru. Kompilator podejmuje decyzje na podstawie struktury konkretnego kodu,
zestawu dostępnych instrukcji procesora oraz przyjętego modelu kosztu. Dodatkowo część flag się wyklucza.
Do udzielenia odpowiedzi potrzebne jest
zatem stanowisko pomiarowe, na którym można uruchomić ten sam kod w~wielu
wariantach kompilacji i~porównać określone metryki optymalizajci.


\section{Cel pracy}
Celem niniejszej pracy jest projekt systemu wbudowanego przeznaczonego do
badania możliwości optymalizacyjnych kompilatora oraz do porównania różnych
metod poprawy wydajności programu. Jako algorytm przetwarzania danych wybrano  szybką
transformatę Fouriera,  algorytm o~ustalonej i~dobrze opisanej strukturze,
szeroko wykorzystywany w~cyfrowym przetwarzaniu sygnałów, a~zarazem na tyle
wymagający obliczeniowo, że różnice w~jakości wygenerowanego kodu są w~jego
przypadku mierzalne.

Na system składają się dwie części. Część sprzętową stanowi
Raspberry Pi 4B, pracująca pod
kontrolą systemu operacyjnego Linux. Część programowa obejmuje jedenaście
wariantów implementacji transformaty Fouriera, różniących się od siebie podejściem do realizacji 
Dyskretnej transformaty Fouriera oraz program
pomiarowy, który uruchamia je w~kontrolowanych warunkach i~rejestruje czas
wykonania wraz z~odczytami niezbędnych informacji o innych metrykach optymalizacji.

Zakres pracy obejmuje trzy grupy zagadnień. Pierwszą z nich jest zbadanie, jaki wpływ
na wydajność ma sama zmiana flag optymalizacji, przy niezmienionym kodzie źródłowym. 
Drugą stanowi porównanie tego efektu z~korzyściami wynikającymi z~celowych modyfikacji programu, takich jak
zmiana układu danych wejściowych, jawne użycie instrukcji wektorowych czy wybór
odmiennego schematu obliczeń. Trzecią jest sprawdzenie, w~jakim stopniu wnioski
te zachowują ważność przy zmianie rozmiaru przetwarzanych danych oraz po
zrównolegleniu obliczeń na wiele rdzeni procesora.
\section{Zawartość rozdziałów}
Ogólną strukturę niniejszej pracy można przedstawić w następujący sposób:
\begin{itemize}

    \item \textbf{Rozdział 2: Architektura systemów operacyjnych w zastosowaniach wbudowanych.}
    \\Omawia rolę systemu operacyjnego w systemie wbudowanym, stosowane w nim abstrakcje i mechanizmy, podział systemów czasu rzeczywistego oraz architektury jądra. Zawiera studium przypadku systemu Linux, użytego jako platforma programowa w niniejszej pracy.

    \item \textbf{Rozdział 3: Algorytm przetwarzania danych i techniki optymalizacji programów komputerowych.}
    \\Przedstawia podstawy matematyczne dyskretnej i szybkiej transformacji Fouriera, przebieg procesu kompilacji oraz wewnętrzną strukturę kompilatora. Zawiera przegląd kompilatorów GCC i Clang, opis przekształceń optymalizujących oraz odpowiadających im flag kompilacji.

    \item \textbf{Rozdział 4: Projekt systemu testowego i metodyka przeprowadzonych badań.}
    \\Opisuje platformę sprzętową, jedenaście zaimplementowanych wariantów transformaty Fouriera oraz architekturę programu pomiarowego. Zawiera badane zestawy flag kompilacji, przyjętą metodykę pomiarów i definicje stosowanych metryk.

    \item \textbf{Rozdział 5: Analiza wyników przeprowadzonych badań.}
    \\Zawiera ocenę powtarzalności pomiarów wraz z progiem istotności różnic, a następnie analizę wyników osobno dla wariantów sekwencyjnych i współbieżnych. Poszczególne podrozdziały omawiają wpływ poziomów optymalizacji, celowych modyfikacji kodu, wyboru kompilatora, rozmiaru danych wejściowych oraz liczby wątków.

    \item \textbf{Rozdział 6: Podsumowanie i wnioski.}
    \\Zawiera wnioski końcowe wynikające z przeprowadzonych badań oraz propozycje kierunków dalszych prac.
\end{itemize}
